\section{Introduction}
\label{sec:introduction}

Writing a research paper is a long sequence of small decisions. A sentence is
shortened, a definition is moved before its first use, a symbol is renamed so
that it no longer collides with another, a proof gains the step a referee
asked for. An agent that edits \LaTeX{} can propose many of these decisions
quickly, but a proposal is only useful when the author can see what it changes
in the typeset paper, decide, and say why. Papilot is built around that loop.

\subsection{Three commitments}

Papilot organises the loop around three commitments.

\begin{enumerate}[label=(C\arabic*)]
  \item \emph{Isolation.} Every task edits its own Git worktree, never the
    manuscript checkout, so that unreviewed text cannot leak into the paper.
  \item \emph{Typeset review.} Changes are reviewed as differences between two
    compiled PDFs, with the \LaTeX{} source one click away.
  \item \emph{Explicit acceptance.} Only accepted changes are merged into the
    working tree, and every decision can be undone.
\end{enumerate}

The first commitment is about safety, the second about legibility and the
third about authority. Together they turn an agent from a co-author that writes
into the manuscript into an assistant that proposes to it.

\subsection{What this paper does}

The rest of this paper makes the commitments precise and studies their
consequences with small, explicit models. Our contributions are:

\begin{itemize}
  \item an algebra of edits and review rounds, with a proof that compatible
    edits commute and can be rebased onto each other (\cref{sec:model}), and
    tight bounds on the number of rounds a task needs;
  \item an estimate of the conflict rate between concurrent tasks under uniform
    and clustered edits (\cref{sec:conflicts});
  \item a content-addressed build scheme in which undoing a decision never
    recompiles (\cref{sec:build});
  \item a geometric account of the PDF view, explaining why a line-wide text
    layer drifts from the printed words and proving that per-word runs bound
    the drift (\cref{sec:geometry});
  \item an algorithm that turns a mark on the page into source lines, and an
    analysis of when its matches are unique (\cref{sec:notes});
  \item queueing and usage models for agents (\cref{sec:agents,sec:usage}),
    the index algebra behind partial commits (\cref{sec:git}) and a security
    model (\cref{sec:security});
  \item a model-based evaluation that tabulates the quantities above
    (\cref{sec:evaluation}).
\end{itemize}

\Cref{sec:using} describes Papilot as an author uses it, \cref{sec:background}
collects the background on \TeX{}, PDF, Git and SyncTeX that the later sections
rely on, \cref{sec:related} discusses related
work, and \cref{sec:conclusion} concludes. Longer proofs are in
\cref{app:proofs}; \cref{app:reference} lists the notation and the keyboard
shortcuts.

\subsection{A note on this document}

This paper is also Papilot's example manuscript. It is deliberately dense with
displayed mathematics, theorems and cross-references, because those are the
places where an agent's edit is easiest to get subtly wrong and a typeset
review most useful. The reader is invited to start a task against any equation
and to review the result.
